% !TEX program = pdflatex
\documentclass[10pt,aspectratio=169]{beamer}

\usepackage[utf8]{inputenc}
\usepackage[T1]{fontenc}
\usepackage[english]{babel}
\usepackage{lmodern}
\usepackage{microtype}
\usepackage{amsmath,amssymb,bm,mathtools}
\usepackage{booktabs,array}

\usetheme[numbering=fraction,progressbar=frametitle]{metropolis}
\setbeamertemplate{frame numbering}[fraction]
\setbeamertemplate{navigation symbols}{}

\definecolor{CourseBlue}{HTML}{1A3C68}
\definecolor{CourseGreen}{HTML}{216E39}
\definecolor{CourseGold}{HTML}{B8860B}
\definecolor{CourseRed}{HTML}{9A2C2C}
\definecolor{SoftBlue}{HTML}{F2F6FA}
\definecolor{SoftGreen}{HTML}{F1F8F3}
\definecolor{SoftRed}{HTML}{FFF3F3}

\setbeamercolor{normal text}{fg=black,bg=white}
\setbeamercolor{structure}{fg=CourseBlue}
\setbeamercolor{title}{fg=CourseBlue}
\setbeamerfont{title}{size=\Large,series=\bfseries}
\setbeamercolor{frametitle}{bg=CourseBlue,fg=white}
\setbeamerfont{frametitle}{series=\bfseries,size=\large}
\setbeamertemplate{frametitle}{
  \nointerlineskip
  \begin{beamercolorbox}[wd=\paperwidth,ht=3.2ex,dp=1.4ex,leftskip=1em,rightskip=1em]{frametitle}
    \usebeamerfont{frametitle}\insertframetitle
  \end{beamercolorbox}
}
\setbeamercolor{block title}{bg=CourseBlue,fg=white}
\setbeamercolor{block body}{bg=SoftBlue,fg=black}
\setbeamercolor{block title example}{bg=CourseGreen,fg=white}
\setbeamercolor{block body example}{bg=SoftGreen,fg=black}
\setbeamercolor{block title alerted}{bg=CourseRed,fg=white}
\setbeamercolor{block body alerted}{bg=SoftRed,fg=black}
\setbeamertemplate{blocks}[rounded][shadow=false]

\newenvironment{modelbox}{\begin{beamercolorbox}[rounded=true,sep=0.7em,wd=\textwidth]{block body}
  \textbf{\color{CourseBlue}MODEL}\quad}{\end{beamercolorbox}}
\newenvironment{marketbox}{\begin{beamercolorbox}[rounded=true,sep=0.7em,wd=\textwidth]{block body example}
  \textbf{\color{CourseGreen}NO-ARBITRAGE}\quad}{\end{beamercolorbox}}
\newenvironment{simplificationbox}{\begin{beamercolorbox}[rounded=true,sep=0.7em,wd=\textwidth]{block body alerted}
  \textbf{\color{CourseRed}WATCH OUT}\quad}{\end{beamercolorbox}}
\newenvironment{pausebox}{\begin{beamercolorbox}[rounded=true,sep=0.8em,wd=\textwidth]{block body}
  \textbf{\color{CourseGold}YOUR TURN}\quad}{\end{beamercolorbox}}

\newcommand{\E}{\mathbb E}
\newcommand{\Q}{\mathbb Q}
\newcommand{\F}{\mathcal F}
\newcommand{\dd}{\,\mathrm d}
\newcommand{\Ncdf}{\Phi}

\AtBeginSection[]{
  \begin{frame}[plain]\vfill
    {\Large\bfseries\color{CourseBlue}\insertsectionhead\par}
    \medskip{\color{CourseBlue}\rule{0.92\textwidth}{0.8pt}}\vfill
  \end{frame}
}

\title{Pricing, Hedging, and Model Risk}
\subtitle{Vanilla claims, path dependence, and stochastic discounting}
\author{}
\date{Lecture 13}

\begin{document}

\begin{frame}[plain]
  \vfill
  {\LARGE\bfseries\color{CourseBlue}Pricing, Hedging, and Model Risk\par}
  \medskip{\large Vanilla claims, path dependence, and stochastic discounting\par}
  \medskip{\color{CourseGold}\rule{0.94\textwidth}{0.8pt}}
  \medskip{Lecture 13\par}\vfill
\end{frame}

\begin{frame}{A price is only half of a model}
\begin{block}{The question}
A model assigns a price, but it also prescribes a hedge. How should we judge
the hedge when the market evolves differently from the model?
\end{block}
\begin{enumerate}
  \item Start with a vanilla call and the Black--Scholes replication benchmark.
  \item Change the model while keeping the contract fixed.
  \item Change the contract while keeping the market model simple.
  \item Add stochastic rates and ask which traded assets span the risks.
\end{enumerate}
\begin{marketbox}
A pricing model should be assessed both by the prices it fits and by the
hedging P\&L it generates outside the model.
\end{marketbox}
\end{frame}

\section{A vanilla option: the hedge depends on the model}

\begin{frame}{For a European call, Black--Scholes links price and replication}
The payoff is
\[
(S_T-K)^+.
\]
Under the risk-neutral Black--Scholes model,
\[
\frac{\dd S_t}{S_t}=r\dd t+\sigma\dd W_t^\Q,
\]
the option price is
\[
u(t,s)=\E^\Q\!\left[e^{-r(T-t)}(S_T-K)^+\mid S_t=s\right].
\]
\begin{block}{The replication target}
A self-financing portfolio of stock and cash must satisfy
\[
\Pi_t=u(t,S_t),\qquad \Pi_T=(S_T-K)^+.
\]
\end{block}
\end{frame}

\begin{frame}{Matching the Brownian exposure determines the delta}
Let \(\Pi_t=\Delta_tS_t+\beta_tB_t\). The self-financing portfolio has
diffusion term
\[
\Delta_t\sigma S_t\dd W_t.
\]
It\^o's formula gives the option diffusion term
\[
\sigma S_tu_s(t,S_t)\dd W_t.
\]
Matching them yields
\[
\boxed{\Delta_t=u_s(t,S_t).}
\]
For a Black--Scholes call,
\[
\boxed{\Delta_t^{\mathrm{BS}}=\Ncdf(d_1(t,S_t)).}
\]
\begin{marketbox}
The PDE enforces the drift; delta neutralizes the random stock exposure.
Pricing and hedging are two consequences of the same model.
\end{marketbox}
\end{frame}

\begin{frame}{A delta hedge must be rebalanced}
For a call, delta moves between \(0\) and \(1\):
\[
\Delta^{\mathrm{BS}}\approx0\quad\text{deep OTM},\qquad
\Delta^{\mathrm{BS}}\approx1\quad\text{deep ITM}.
\]
\begin{pausebox}
A call has \(S_0=100\), \(K=100\), \(r=0\), \(\sigma=20\%\), and \(T=1\).
At inception, \(d_1=0.10\), so \(\Delta_0=\Phi(0.10)\approx0.54\).
What position neutralizes the stock exposure of a short call?
\end{pausebox}
\pause
\begin{exampleblock}{Answer}
Hold approximately \(0.54\) shares per short call, and finance the residual
position with cash. Rebalance after the stock, time, or model inputs change.
\end{exampleblock}
\end{frame}

\begin{frame}{The same call has a different delta under local volatility}
Under local volatility,
\[
\frac{\dd S_t}{S_t}=r\dd t+\sigma_{\mathrm{loc}}(t,S_t)\dd W_t^\Q.
\]
The price still solves a one-factor PDE, but with
\[
\tfrac12\sigma^2s^2u_{ss}
\quad\longrightarrow\quad
\tfrac12\sigma_{\mathrm{loc}}(t,s)^2s^2u_{ss}.
\]
The hedge remains formally
\[
\boxed{\Delta_t^{\mathrm{LV}}=u_s^{\mathrm{LV}}(t,S_t),}
\]
but generally
\[
\Delta_t^{\mathrm{LV}}\neq\Delta_t^{\mathrm{BS}}.
\]
\begin{simplificationbox}
The payoff did not change. The price function changed, so the hedge changed.
A delta is never model-free.
\end{simplificationbox}
\end{frame}

\begin{frame}{A market quote does not eliminate the need for a model}
A liquid vanilla call already has an observable market price. Why use a local
volatility model for it?
\begin{columns}[T,onlytextwidth]
\column{0.48\textwidth}
\begin{block}{What the market quote gives}
Today's value of the option.
\end{block}
\column{0.48\textwidth}
\begin{exampleblock}{What the model gives}
A dynamic law, a price surface between quotes, and a hedge consistent with
the observed smile.
\end{exampleblock}
\end{columns}
\begin{marketbox}
Local volatility calibrates the current vanilla surface; its delta reflects
the smile information ignored by a flat-volatility Black--Scholes delta.
\end{marketbox}
\end{frame}

\begin{frame}{Hedging P\&L is the test performed outside the model}
Suppose an option is sold and hedged dynamically. At maturity,
\[
\boxed{\mathrm{P\&L}=\Pi_T-(S_T-K)^+.}
\]
\begin{center}
\begin{tabular}{c|l}
\toprule
Source of nonzero P\&L & Why it remains\\
\midrule
Discrete trading & Rebalancing is not continuous\\
Numerics and costs & Greeks are approximated; trading is not frictionless\\
Model mismatch & The market does not follow the hedging model\\
\bottomrule
\end{tabular}
\end{center}
\begin{pausebox}
Does a nonzero terminal P\&L prove that the delta formula was wrong?
\end{pausebox}
\pause
\begin{exampleblock}{Answer}
No. Discrete trading and model mismatch create P\&L even with a correct
model delta.
\end{exampleblock}
\end{frame}

\begin{frame}{Model risk becomes visible under stochastic volatility}
Imagine the market follows a Heston-type model,
\[
\frac{\dd S_t}{S_t}=r\dd t+\sqrt{V_t}\dd W_t^{(1)},\qquad
\dd V_t=\kappa(\theta-V_t)\dd t+\eta\sqrt{V_t}\dd W_t^{(2)}.
\]
A Black--Scholes hedge ignores both smile and random variance. A
local-volatility hedge incorporates today's smile, but still makes variance
a deterministic function of time and spot.
\begin{block}{What a simulation experiment compares}
Simulate the same market paths, rebalance each hedge, then compare the
distributions of their terminal P\&L: mean, dispersion, and tail losses.
\end{block}
\begin{marketbox}
A more market-consistent model may improve hedging, but no model dominates
for every contract, parameter set, or rebalancing frequency.
\end{marketbox}
\end{frame}

\section{A lookback option: the path becomes a state variable}

\begin{frame}{A lookback call remembers the highest stock level}
A fixed-strike lookback call pays
\[
\boxed{(M_T-K)^+,\qquad M_T=\max_{0\leq u\leq T}S_u.}
\]
\begin{columns}[T,onlytextwidth]
\column{0.48\textwidth}
\begin{block}{Vanilla call}
\[
(S_T-K)^+.
\]
Only the terminal stock price matters.
\end{block}
\column{0.48\textwidth}
\begin{exampleblock}{Lookback call}
\[
(M_T-K)^+.
\]
The whole price path matters through its running maximum.
\end{exampleblock}
\end{columns}
\end{frame}

\begin{frame}{Same terminal stock, different path, different payoff}
Take \(K=100\). Two paths both finish at \(S_T=100\):
\[
\begin{array}{c|c|c}
\toprule
\text{Path} & M_T & \text{Lookback payoff}\\
\midrule
A & 120 & (120-100)^+=20\\
B & 105 & (105-100)^+=5\\
\bottomrule
\end{array}
\]
\begin{pausebox}
Which piece of information is missing if we know only \(S_T=100\)?
\end{pausebox}
\pause
\begin{exampleblock}{Answer}
The running maximum. The current or terminal spot alone cannot summarize a
path-dependent payoff.
\end{exampleblock}
\end{frame}

\begin{frame}{State augmentation makes the path-dependent problem Markov}
At time \(t\), define
\[
M_t=\max_{0\leq u\leq t}S_u.
\]
The correct state is
\[
\boxed{(t,S_t,M_t),\qquad S_t\leq M_t.}
\]
The pricing function becomes
\[
\boxed{
u(t,s,m)=\E^\Q\!\left[
e^{-r(T-t)}(M_T-K)^+\mid S_t=s,\ M_t=m
\right].}
\]
\begin{marketbox}
The full past is not stored. We retain only the smallest path statistic that
matters for the future payoff.
\end{marketbox}
\end{frame}

\begin{frame}{The running maximum changes only at a new record}
\begin{center}
\begin{tabular}{c|c|c}
\toprule
Current state & Next stock move & What happens to \(M\)?\\
\midrule
\(S_t<M_t\) & small move below \(M_t\) & \(M\) stays fixed\\
\(S_t=M_t\) & stock moves down & \(M\) stays fixed\\
\(S_t=M_t\) & stock moves up & \(M\) increases with the stock\\
\bottomrule
\end{tabular}
\end{center}
\begin{block}{Why this complicates a PDE}
The state domain is \(s\leq m\), and the diagonal \(s=m\) carries a special
boundary condition. The extra variable is not an independent Brownian factor;
it is a constrained path functional.
\end{block}
\end{frame}

\begin{frame}{Path dependence makes model choice more consequential}
\begin{columns}[T,onlytextwidth]
\column{0.48\textwidth}
\begin{block}{Vanilla option}
The payoff depends mainly on the terminal distribution of \(S_T\). Liquid
vanillas constrain that distribution directly.
\end{block}
\column{0.48\textwidth}
\begin{alertblock}{Lookback option}
The payoff depends on hitting high levels along the way. Different models can
fit vanillas today yet assign different probabilities to those path events.
\end{alertblock}
\end{columns}
\begin{marketbox}
For exotics, model risk is not merely a question of the current smile. It is
a question about future path dynamics.
\end{marketbox}
\end{frame}

\begin{frame}{Pricing a lookback option requires richer numerical state}
\begin{columns}[T,onlytextwidth]
\column{0.48\textwidth}
\begin{block}{Monte Carlo}
Simulate \(S\) path by path, update
\[
M_{t_{n+1}}=\max\{M_{t_n},S_{t_{n+1}}\},
\]
then discount the terminal payoff.
\end{block}
\column{0.48\textwidth}
\begin{exampleblock}{PDE approach}
Solve for \(u(t,s,m)\) on the constrained domain \(s\leq m\), with a
special condition on the record boundary \(s=m\).
\end{exampleblock}
\end{columns}
\begin{simplificationbox}
A monitoring grid can miss extrema between dates. For a lookback, this
discretization choice directly affects the payoff and hedge.
\end{simplificationbox}
\end{frame}

\begin{frame}{The lookback delta depends on history as well as spot}
The hedge is still a sensitivity to the tradable stock:
\[
\Delta_t^{\mathrm{LB}}=u_s(t,S_t,M_t).
\]
\begin{pausebox}
Two contracts have the same \(t\) and \(S_t=105\). Contract A has
\(M_t=120\), contract B has \(M_t=108\). Should their deltas generally be
identical?
\end{pausebox}
\pause
\begin{exampleblock}{Answer}
No. The established maximum changes the remaining payoff opportunity. The
delta must reflect both spot and the historical record \(M_t\).
\end{exampleblock}
\end{frame}

\begin{frame}{Exotics turn model error into hedging error faster}
For a lookback hedge, terminal P\&L reflects the usual sources of error plus
path-specific ones:
\[
\begin{array}{c|l}
\toprule
\text{Risk} & \text{Effect}\\
\midrule
Smile dynamics & changes probabilities of future records\\
Volatility dynamics & changes the frequency of extreme moves\\
Monitoring frequency & can miss a short-lived maximum\\
Delta interpolation & is needed on the enlarged state grid\\
\bottomrule
\end{array}
\]
\begin{marketbox}
The model used for an exotic should be stress-tested against the scenarios
that drive its path functional, not only against vanilla prices.
\end{marketbox}
\end{frame}

\section{Stochastic rates add a hedging factor}

\begin{frame}{Random rates make discounting and the state variable random}
Let \(r_t\) be a stochastic short rate. Even for a payoff \(g(S_T)\),
\[
\boxed{
u(t,s,r)=\E^\Q\!\left[
e^{-\int_t^T r_u\dd u}g(S_T)\mid S_t=s,\ r_t=r
\right].}
\]
\begin{block}{What changed?}
The state is now \((t,S_t,r_t)\). Rates affect the discount factor, the
risk-neutral drift of the stock, and potentially the correlation structure.
\end{block}
\end{frame}

\begin{frame}{Stock delta alone does not span an independent rate shock}
Suppose the market has two Brownian sources: one stock-related and one
rate-related.
\begin{center}
\begin{tabular}{c|c|c}
\toprule
Traded position & Stock risk & Rate risk\\
\midrule
Stock & yes & possibly partial\\
Cash account & no diffusion exposure & no diffusion exposure\\
Zero-coupon bond & rate sensitivity & yes\\
\bottomrule
\end{tabular}
\end{center}
\begin{pausebox}
Can stock plus cash perfectly hedge an option exposed to an independent rate
factor?
\end{pausebox}
\pause
\begin{alertblock}{No}
Cash does not add a new Brownian exposure. A rate-sensitive risky instrument,
such as a bond or swap, is needed to hedge an independent rate factor.
\end{alertblock}
\end{frame}

\begin{frame}{The natural two-factor hedge is stock plus bond plus cash}
Choose a zero-coupon bond \(P(t,T^\star)\) and form
\[
\Pi_t=\phi_tS_t+\psi_tP(t,T^\star)+\beta_tB_t.
\]
If stock and bond exposures span the two Brownian directions, choose
\(\phi_t,\psi_t\) to solve
\[
\begin{pmatrix}
\Gamma_S^{(1)} & \Gamma_P^{(1)}\\
\Gamma_S^{(2)} & \Gamma_P^{(2)}
\end{pmatrix}
\begin{pmatrix}\phi_t\\ \psi_t\end{pmatrix}
=
\begin{pmatrix}\Gamma_u^{(1)}\\ \Gamma_u^{(2)}\end{pmatrix}.
\]
\begin{marketbox}
Exact local replication is possible when this exposure matrix has full rank.
The cash account then matches the value and drift.
\end{marketbox}
\end{frame}

\begin{frame}{Stochastic rates do not automatically imply incompleteness}
\begin{columns}[T,onlytextwidth]
\column{0.48\textwidth}
\begin{exampleblock}{Complete case}
Every relevant Brownian risk is spanned by traded assets, for example stock
and a sufficiently sensitive bond in a two-factor setting.
\end{exampleblock}
\column{0.48\textwidth}
\begin{alertblock}{Incomplete case}
There are more risk factors than independent traded exposures, for example
stock plus cash but an independent short-rate factor.
\end{alertblock}
\end{columns}
\begin{pausebox}
What determines completeness: whether rates are random, or whether tradable
assets span the risk factors?
\end{pausebox}
\pause
\begin{exampleblock}{Answer}
Spanning. Random rates enlarge the hedging problem, but they do not by
themselves prevent exact replication.
\end{exampleblock}
\end{frame}

\begin{frame}{In an incomplete market, hedging is risk management}
When perfect replication is unavailable, the objective becomes:
\begin{enumerate}
  \item hedge stock sensitivity with the stock;
  \item hedge rate sensitivity with bonds, swaps, or futures;
  \item monitor and limit the residual, unspanned P\&L.
\end{enumerate}
\begin{block}{A useful distinction}
\[
\text{complete market: replication determines price}
\]
\[
\text{incomplete market: model and risk criterion select a hedge}
\]
\end{block}
\begin{marketbox}
The complete-market benchmark remains valuable: it identifies which missing
instruments or risk factors create the residual error.
\end{marketbox}
\end{frame}

\section{Takeaways}

\begin{frame}{Three case studies, one common lesson}
\begin{center}
\begin{tabular}{p{0.24\textwidth}|p{0.30\textwidth}|p{0.29\textwidth}}
\toprule
Case & New source of complexity & Consequence for hedging\\
\midrule
Vanilla call & model choice & delta is model dependent\\
Lookback option & path dependence & state and delta include the running maximum\\
Stochastic rates & extra rate factor & stock hedge must be complemented by bonds\\
\bottomrule
\end{tabular}
\end{center}
\begin{marketbox}
As contracts or models add state variables, the hedge must account for the
new sources of risk or accept residual P\&L.
\end{marketbox}
\end{frame}

\begin{frame}{Final check: judge the hedge where it will be used}
\begin{pausebox}
A desk calibrates a model perfectly to today's vanilla smile, then hedges a
one-year lookback option daily. Which additional questions should it ask
before trusting a small model price error?
\end{pausebox}
\pause
\begin{exampleblock}{Answer}
Does the model generate plausible future smile and path dynamics? How large
are P\&L tails under alternative models and monitoring grids? Can the
required Greeks be hedged with liquid instruments?
\end{exampleblock}
\vfill
\begin{center}
{\large\color{CourseBlue}\bfseries A model price is a starting point; hedging P\&L is the operational test.}
\end{center}
\end{frame}

\end{document}
